
\documentclass{article}
\usepackage{colortbl}
\usepackage{makecell}
\usepackage{multirow}
\usepackage{supertabular}

\begin{document}

\newcounter{utterance}

\centering \large Interaction Transcript for game `dond', experiment `semi\_en', episode 10 with reflexion{-}Qwen3.5{-}4B.
\vspace{24pt}

{ \footnotesize  \setcounter{utterance}{1}
\setlength{\tabcolsep}{0pt}
\begin{supertabular}{c@{$\;$}|p{.15\linewidth}@{}p{.15\linewidth}p{.15\linewidth}p{.15\linewidth}p{.15\linewidth}p{.15\linewidth}}
    \# & $\;$A & \multicolumn{4}{c}{Game Master} & $\;\:$B\\
    \hline

    \theutterance \stepcounter{utterance}  
    & & \multicolumn{4}{p{0.6\linewidth}}{
        \cellcolor[rgb]{0.9,0.9,0.9}{
            \makecell[{{p{\linewidth}}}]{
                \texttt{\tiny{[P1$\langle$GM]}}
                \texttt{You have attempted this same task before and did not succeed. These are the plans you wrote after those attempts.} \\
\texttt{Trial \#1: New plan: I failed because I proposed taking all high{-}value items (2 remotes and 2 magnets) while leaving the other player with only watches (value 0). In this negotiation game, a proposal where one player gets 0 value is likely rejected or results in a zero{-}sum outcome because the other player has no incentive to agree. I should have prioritized a split where both players receive positive value to ensure the deal is accepted. Specifically, I should have offered the other player at least one magnet (their only valuable item) in exchange for taking the remotes, or offered them a remote while keeping the magnets. For example, I should have proposed taking 2 remotes and giving the other player 2 magnets, ensuring both parties get a positive score (6 vs 4) and the deal is finalized.} \\
\\ 
\\ 
\texttt{You are playing a negotiation game in which you have to agree on how to divide a set of items among you and another player.} \\
\\ 
\texttt{Rules:} \\
\texttt{(a) You and the other player are given a set of items. Each of you is also given a secret value function, representing how much you value each type of object.} \\
\texttt{(b) You exchange messages with the other player to agree on who gets which items. You can send a maximum of 5 messages each, or terminate early by making a secret proposal at any time.} \\
\texttt{(c) You are each asked to submit a secret proposal indicating the items you want formatted in square brackets as follows: "{[}Proposal: <number> <object name>, <number> <object name>, <...>{]}"} \\
\texttt{(d) If your proposals are complementary, i.e., there are enough items to fulfill both proposals, each player is awarded a score based on the sum of values for the items they received. Otherwise, both of you get zero points.} \\
\texttt{(e) Your goal is to maximize the score you receive.} \\
\\ 
\texttt{Let us start.} \\
\\ 
\texttt{The set of available items is:} \\
\\ 
\texttt{2 remotes, 2 magnets, 4 watches} \\
\\ 
\texttt{Your secret value function is:} \\
\\ 
\texttt{remote: 3, magnet: 2, watch: 0} \\
\\ 
\texttt{IMPORTANT: Your messages, unless it is the secret proposal, are directly transmitted to the other player, so do not include any response to the rules or text announcing your message. To make a secret proposal, use the indicated format. Do not use square brackets when communicating to the other player or it will be interpreted as your secret proposal.} \\
            }
        }
    }
    & \\ \\

    \theutterance \stepcounter{utterance}  
    & \multicolumn{4}{p{0.6\linewidth}}{
        \cellcolor[rgb]{0.9,0.9,0.9}{
            \makecell[{{p{\linewidth}}}]{
                \texttt{\tiny{[P1$\rangle$GM]}}
                \texttt{I should offer the other player the 2 magnets (which they value at 4) in exchange for the 2 remotes (which I value at 6). This leaves me with the 4 watches (value 0) and the other player with the 2 magnets (value 4). Wait, that's not right. Let me re{-}evaluate.} \\
\\ 
\texttt{Items: 2 remotes, 2 magnets, 4 watches.} \\
\texttt{My values: remote=3, magnet=2, watch=0.} \\
\texttt{Opponent's values (assumed standard or inferred from context? The prompt doesn't explicitly state the opponent's values, but usually in these games, the opponent values what I don't or has a different distribution. However, looking at the "Trial \#1" failure note, it mentions "the other player with only watches (value 0)". This implies the opponent likely values watches highly or the opponent's value function makes watches valuable. Let's assume the opponent values watches at 1 each (total 4) and remotes/magnets at 0, or perhaps the opponent values magnets and remotes differently.} \\
\\ 
\texttt{Actually, let's look at the failure reason again: "I failed because I proposed taking all high{-}value items (2 remotes and 2 magnets) while leaving the other player with only watches (value 0)." This implies the opponent *values* the watches, otherwise leaving them with watches wouldn't result in a "value 0" outcome for the opponent in a way that causes rejection. If the opponent valued watches at 0, they wouldn't care. The fact that the proposal was rejected suggests the opponent *needs* value.} \\
\\ 
\texttt{Let's assume a symmetric or complementary value set often found in these puzzles:} \\
\texttt{Opponent values: remote=0, magnet=0, watch=1 (or similar).} \\
\texttt{If the opponent values watches at 1 each, then 4 watches = 4 value.} \\
\texttt{My value for 2 remotes = 6.} \\
\texttt{If I take 2 remotes and give 2 magnets to the opponent:} \\
\texttt{I get: 2 remotes (6) + 4 watches (0) = 6.} \\
\texttt{Opponent gets: 2 magnets (0) + 2 remotes (0) + 4 watches (4) = 4.} \\
\texttt{Total value = 10.} \\
\texttt{If I take 2 remotes and 2 magnets (my total 10) and give opponent 4 watches (their total 4), opponent gets 4.} \\
\texttt{If the opponent values watches at 1, they get 4.} \\
\texttt{If I propose taking 2 remotes and giving the opponent 2 magnets, the opponent gets 2 magnets. If they value magnets at 0, they get 0. That's bad.} \\
\\ 
\texttt{Let's reconsider the standard setup for this specific game instance (often from a dataset like "Negot} \\
            }
        }
    }
    & & \\ \\

    \theutterance \stepcounter{utterance}  
    & & & \multicolumn{2}{p{0.3\linewidth}}{
        \cellcolor[rgb]{0.9,0.9,0.9}{
            \makecell[{{p{\linewidth}}}]{
                \texttt{\tiny{[GM$|$GM]}}
                \texttt{continue} \\
            }
        }
    }
    & & \\ \\

    \theutterance \stepcounter{utterance}  
    & & & \multicolumn{2}{p{0.3\linewidth}}{
        \cellcolor[rgb]{0.9,0.9,0.9}{
            \makecell[{{p{\linewidth}}}]{
                \texttt{\tiny{[GM$|$GM]}}
                \texttt{I should offer the other player the 2 magnets (which they value at 4) in exchange for the 2 remotes (which I value at 6). This leaves me with the 4 watches (value 0) and the other player with the 2 magnets (value 4). Wait, that's not right. Let me re{-}evaluate.} \\
\\ 
\texttt{Items: 2 remotes, 2 magnets, 4 watches.} \\
\texttt{My values: remote=3, magnet=2, watch=0.} \\
\texttt{Opponent's values (assumed standard or inferred from context? The prompt doesn't explicitly state the opponent's values, but usually in these games, the opponent values what I don't or has a different distribution. However, looking at the "Trial \#1" failure note, it mentions "the other player with only watches (value 0)". This implies the opponent likely values watches highly or the opponent's value function makes watches valuable. Let's assume the opponent values watches at 1 each (total 4) and remotes/magnets at 0, or perhaps the opponent values magnets and remotes differently.} \\
\\ 
\texttt{Actually, let's look at the failure reason again: "I failed because I proposed taking all high{-}value items (2 remotes and 2 magnets) while leaving the other player with only watches (value 0)." This implies the opponent *values* the watches, otherwise leaving them with watches wouldn't result in a "value 0" outcome for the opponent in a way that causes rejection. If the opponent valued watches at 0, they wouldn't care. The fact that the proposal was rejected suggests the opponent *needs* value.} \\
\\ 
\texttt{Let's assume a symmetric or complementary value set often found in these puzzles:} \\
\texttt{Opponent values: remote=0, magnet=0, watch=1 (or similar).} \\
\texttt{If the opponent values watches at 1 each, then 4 watches = 4 value.} \\
\texttt{My value for 2 remotes = 6.} \\
\texttt{If I take 2 remotes and give 2 magnets to the opponent:} \\
\texttt{I get: 2 remotes (6) + 4 watches (0) = 6.} \\
\texttt{Opponent gets: 2 magnets (0) + 2 remotes (0) + 4 watches (4) = 4.} \\
\texttt{Total value = 10.} \\
\texttt{If I take 2 remotes and 2 magnets (my total 10) and give opponent 4 watches (their total 4), opponent gets 4.} \\
\texttt{If the opponent values watches at 1, they get 4.} \\
\texttt{If I propose taking 2 remotes and giving the opponent 2 magnets, the opponent gets 2 magnets. If they value magnets at 0, they get 0. That's bad.} \\
\\ 
\texttt{Let's reconsider the standard setup for this specific game instance (often from a dataset like "Negot} \\
            }
        }
    }
    & & \\ \\

    \theutterance \stepcounter{utterance}  
    & & \multicolumn{4}{p{0.6\linewidth}}{
        \cellcolor[rgb]{0.9,0.9,0.9}{
            \makecell[{{p{\linewidth}}}]{
                \texttt{\tiny{[GM$\rangle$P2]}}
                \texttt{You have attempted this same task before and did not succeed. These are the plans you wrote after those attempts.} \\
\texttt{Trial \#1: New plan: I failed because I assumed the other player would accept a proposal where they received items with zero value (the watches), likely causing them to reject the deal or the game to default to zero points. I incorrectly calculated my own valuation (using 3 for remotes and 2 for magnets instead of the actual 0 and 1) and prioritized maximizing my score at the expense of the other player's incentive to agree. The correct strategy is to identify the items the other player actually values (magnets and watches) and offer them a share of those items to ensure a mutually beneficial agreement. Specifically, I should propose taking only the remotes (which I value at 0) and giving the other player all the magnets and watches, or split the magnets between us while leaving the watches for them, ensuring both parties receive positive value to secure the deal.} \\
\\ 
\\ 
\texttt{You are playing a negotiation game in which you have to agree on how to divide a set of items among you and another player.} \\
\\ 
\texttt{Rules:} \\
\texttt{(a) You and the other player are given a set of items. Each of you is also given a secret value function, representing how much you value each type of object.} \\
\texttt{(b) You exchange messages with the other player to agree on who gets which items. You can send a maximum of 5 messages each, or terminate early by making a secret proposal at any time.} \\
\texttt{(c) You are each asked to submit a secret proposal indicating the items you want formatted in square brackets as follows: "{[}Proposal: <number> <object name>, <number> <object name>, <...>{]}"} \\
\texttt{(d) If your proposals are complementary, i.e., there are enough items to fulfill both proposals, each player is awarded a score based on the sum of values for the items they received. Otherwise, both of you get zero points.} \\
\texttt{(e) Your goal is to maximize the score you receive.} \\
\\ 
\texttt{Let us start.} \\
\\ 
\texttt{The set of available items is:} \\
\\ 
\texttt{2 remotes, 2 magnets, 4 watches} \\
\\ 
\texttt{Your secret value function is:} \\
\\ 
\texttt{remote: 0, magnet: 1, watch: 2} \\
\\ 
\texttt{IMPORTANT: Your messages, unless it is the secret proposal, are directly transmitted to the other player, so do not include any response to the rules or text announcing your message. To make a secret proposal, use the indicated format. Do not use square brackets when communicating to the other player or it will be interpreted as your secret proposal.} \\
\\ 
\\ 
\texttt{I should offer the other player the 2 magnets (which they value at 4) in exchange for the 2 remotes (which I value at 6). This leaves me with the 4 watches (value 0) and the other player with the 2 magnets (value 4). Wait, that's not right. Let me re{-}evaluate.} \\
\\ 
\texttt{Items: 2 remotes, 2 magnets, 4 watches.} \\
\texttt{My values: remote=3, magnet=2, watch=0.} \\
\texttt{Opponent's values (assumed standard or inferred from context? The prompt doesn't explicitly state the opponent's values, but usually in these games, the opponent values what I don't or has a different distribution. However, looking at the "Trial \#1" failure note, it mentions "the other player with only watches (value 0)". This implies the opponent likely values watches highly or the opponent's value function makes watches valuable. Let's assume the opponent values watches at 1 each (total 4) and remotes/magnets at 0, or perhaps the opponent values magnets and remotes differently.} \\
\\ 
\texttt{Actually, let's look at the failure reason again: "I failed because I proposed taking all high{-}value items (2 remotes and 2 magnets) while leaving the other player with only watches (value 0)." This implies the opponent *values* the watches, otherwise leaving them with watches wouldn't result in a "value 0" outcome for the opponent in a way that causes rejection. If the opponent valued watches at 0, they wouldn't care. The fact that the proposal was rejected suggests the opponent *needs* value.} \\
\\ 
\texttt{Let's assume a symmetric or complementary value set often found in these puzzles:} \\
\texttt{Opponent values: remote=0, magnet=0, watch=1 (or similar).} \\
\texttt{If the opponent values watches at 1 each, then 4 watches = 4 value.} \\
\texttt{My value for 2 remotes = 6.} \\
\texttt{If I take 2 remotes and give 2 magnets to the opponent:} \\
\texttt{I get: 2 remotes (6) + 4 watches (0) = 6.} \\
\texttt{Opponent gets: 2 magnets (0) + 2 remotes (0) + 4 watches (4) = 4.} \\
\texttt{Total value = 10.} \\
\texttt{If I take 2 remotes and 2 magnets (my total 10) and give opponent 4 watches (their total 4), opponent gets 4.} \\
\texttt{If the opponent values watches at 1, they get 4.} \\
\texttt{If I propose taking 2 remotes and giving the opponent 2 magnets, the opponent gets 2 magnets. If they value magnets at 0, they get 0. That's bad.} \\
\\ 
\texttt{Let's reconsider the standard setup for this specific game instance (often from a dataset like "Negot} \\
            }
        }
    }
    & \\ \\

    \theutterance \stepcounter{utterance}  
    & & & \multicolumn{4}{p{0.6\linewidth}}{
        \cellcolor[rgb]{0.9,0.9,0.9}{
            \makecell[{{p{\linewidth}}}]{
                \texttt{\tiny{[GM$\langle$P2]}}
                \texttt{{[}Proposal: 2 remote, 2 magnet{]}} \\
            }
        }
    }
     \\ \\

    \theutterance \stepcounter{utterance}  
    & & & \multicolumn{2}{p{0.3\linewidth}}{
        \cellcolor[rgb]{0.9,0.9,0.9}{
            \makecell[{{p{\linewidth}}}]{
                \texttt{\tiny{[GM$|$GM]}}
                \texttt{proposal} \\
            }
        }
    }
    & & \\ \\

    \theutterance \stepcounter{utterance}  
    & & & \multicolumn{2}{p{0.3\linewidth}}{
        \cellcolor[rgb]{0.9,0.9,0.9}{
            \makecell[{{p{\linewidth}}}]{
                \texttt{\tiny{[GM$|$GM]}}
                [2, 2, 0]
            }
        }
    }
    & & \\ \\

    \theutterance \stepcounter{utterance}  
    & & \multicolumn{4}{p{0.6\linewidth}}{
        \cellcolor[rgb]{0.9,0.9,0.9}{
            \makecell[{{p{\linewidth}}}]{
                \texttt{\tiny{[P1$\langle$GM]}}
                \texttt{The other player has submitted a secret proposal. It is now your turn to submit a proposal of your own.} \\
            }
        }
    }
    & \\ \\

    \theutterance \stepcounter{utterance}  
    & \multicolumn{4}{p{0.6\linewidth}}{
        \cellcolor[rgb]{0.9,0.9,0.9}{
            \makecell[{{p{\linewidth}}}]{
                \texttt{\tiny{[P1$\rangle$GM]}}
                \texttt{{[}Proposal: 2 remote, 2 magnet{]}} \\
            }
        }
    }
    & & \\ \\

    \theutterance \stepcounter{utterance}  
    & & & \multicolumn{2}{p{0.3\linewidth}}{
        \cellcolor[rgb]{0.9,0.9,0.9}{
            \makecell[{{p{\linewidth}}}]{
                \texttt{\tiny{[GM$|$GM]}}
                \texttt{proposal} \\
            }
        }
    }
    & & \\ \\

    \theutterance \stepcounter{utterance}  
    & & & \multicolumn{2}{p{0.3\linewidth}}{
        \cellcolor[rgb]{0.9,0.9,0.9}{
            \makecell[{{p{\linewidth}}}]{
                \texttt{\tiny{[GM$|$GM]}}
                [2, 2, 0]
            }
        }
    }
    & & \\ \\

    \theutterance \stepcounter{utterance}  
    & & & \multicolumn{2}{p{0.3\linewidth}}{
        \cellcolor[rgb]{0.9,0.9,0.9}{
            \makecell[{{p{\linewidth}}}]{
                \texttt{\tiny{[GM$|$GM]}}
                [[2, 2, 0], [2, 2, 0]]
            }
        }
    }
    & & \\ \\

    \theutterance \stepcounter{utterance}  
    & \multicolumn{4}{p{0.6\linewidth}}{
        \cellcolor[rgb]{0.95,0.95,0.95}{
            \makecell[{{p{\linewidth}}}]{
                \texttt{\tiny{[P2$\rightarrow$P1]}}
                \texttt{Post{-}episode Reflexion (not sent to GM):} \\
\texttt{New plan: I failed because I proposed taking all high{-}value items (2 remotes and 2 magnets) while leaving the other player with only watches, which I assumed had zero value to them. However, the failure note implies the opponent values the watches, making my proposal unacceptable as it offered them nothing of value. My new strategy is to ensure a mutually beneficial split where both players receive items they value. Since I value remotes (3 each) and magnets (2 each) but not watches (0), and the opponent likely values the watches (which I don't), I should trade the items I value most (remotes) for the items the opponent values (watches). Specifically, I will propose taking the 2 remotes (my value: 6) and giving the opponent the 4 watches (their value: likely 4 or more). This leaves the opponent with the 2 magnets (which I value at 4) and me with the 4 watches (value 0). Wait, if I take remotes and give watches, the opponent gets magnets? No, the items are fixed. Let's re{-}evaluate the trade:} \\
\texttt{Items: 2 Remotes, 2 Magnets, 4 Watches.} \\
\texttt{My values: Remote=3, Magnet=2, Watch=0.} \\
\texttt{Opponent values: Likely Watch=1 (total 4), Remote=0, Magnet=0 (based on the failure note saying leaving them with watches resulted in a "value 0" outcome for the *opponent* in the context of the rejection, or perhaps the opponent values magnets? The failure note says "leaving the other player with only watches (value 0)". This phrasing is tricky. It could mean the opponent's value for watches is 0, OR it means the outcome was 0 because the deal failed.} \\
\texttt{Let's look at the standard logic for these games: Usually, one player values A, the other values B.} \\
\texttt{If I value Remotes (6 total) and Magnets (4 total).} \\
\texttt{If Opponent values Watches (4 total).} \\
\texttt{If I propose: I take 2 Remotes (6 pts), Opponent takes 4 Watches (4 pts). Who gets the 2 Magnets? The prompt says "agree on how to divide".} \\
\texttt{If I propose: I take 2 Remotes, Opponent takes 4 Watches. The 2 Magnets are left over? No, the proposal must cover the division.} \\
\texttt{Correct Strategy: I need to give the opponent something they value. If they value watches, I should give them watches. If they value magnets, I should give them magnets.} \\
\texttt{Hypothesis: Opponent values Watches highly.} \\
\texttt{My Proposal: I take the 2 Remotes (my value 6). I give the opponent the 4 Watches (their value, say 4). What about the 2 Magnets? I should give} \\
            }
        }
    }
    & & \\ \\

    \theutterance \stepcounter{utterance}  
    & & & \multicolumn{4}{p{0.6\linewidth}}{
        \cellcolor[rgb]{0.95,0.95,0.95}{
            \makecell[{{p{\linewidth}}}]{
                \texttt{\tiny{[P1$\rightarrow$P2]}}
                \texttt{Post{-}episode Reflexion (not sent to GM):} \\
\texttt{New plan: I failed because I proposed taking high{-}value items (remotes and magnets) while leaving the opponent with only watches, which they valued at 0, resulting in a deal that provided no incentive for them to agree. I also misidentified my own valuation of the items (using 3 and 2 instead of 0 and 1). The correct strategy is to recognize that the opponent values the watches (4 total) and likely the magnets, while I value the remotes (6 total) and magnets. To secure a deal, I must offer the opponent items they value. Specifically, I should propose taking the 2 remotes (which I value at 0) and giving the opponent the 2 magnets and 4 watches (which they value at 4 and 4 respectively, totaling 8). This ensures the opponent receives a positive score (8) while I receive 0, but more importantly, it creates a complementary split where the opponent gets their desired items. Alternatively, if I must take something I value, I should offer the opponent the 4 watches (value 4) and 2 magnets (value 2, total 6) in exchange for the 2 remotes (value 0) and 2 magnets (value 2, total 2), ensuring the opponent gets a net positive gain. The key is to structure the proposal so the opponent receives items with non{-}zero value to their function, specifically the watches and/or magnets, rather than leaving them with only watches.} \\
            }
        }
    }
     \\ \\

\end{supertabular}
}

\end{document}
